\documentclass[10pt,a4paper]{amsart}
\usepackage[margin=19mm]{geometry}
\usepackage{amsmath,amssymb,mathtools}
\usepackage[scr=rsfs,cal=euler]{mathalfa}
\usepackage{xfrac}
\usepackage[unicode,hidelinks,bookmarks=false]{hyperref}
\newcommand{\ie}{i.e.\ }
\newcommand{\cf}{cf.\ }
\newcommand{\resp}{respectively\ }
\newcommand{\Subsection}{\subsection}
\providecommand{\nobreakdash}{\nobreak\hbox{-}}
\setlength{\emergencystretch}{2em}
\multlinegap0pt
\usepackage{amssymb}
\usepackage{cancel}
\usepackage{tikz}
\usepackage{float}
\usepackage[shortlabels]{enumitem}
\newtheorem{assumption}{Assumption}
\newtheorem{proposition}{Proposition}
\newtheorem{theorem}{Theorem}
\renewcommand{\eqref}[1]{Eq.~(\ref{#1})}  %Modified equation reference
\newcommand{\vect}[1]{\boldsymbol{#1}}
\title{\LARGE \bf Modelling the coevolution of opinion dynamics and decision making in social dilemmas}
\author{Ella C. Davidson, Lorenzo Zino, Ming Cao, Mengbin Ye}
\date{Source: arXiv:2604.08840v1 (CC BY 4.0)}
\begin{document}
\maketitle

\noindent\emph{Source and licence.} \LARGE \bf Modelling the coevolution of opinion dynamics and decision making in social dilemmas. Version arXiv:2604.08840v1 (CC BY 4.0), distributed by arXiv under the Creative Commons Attribution 4.0 International licence, http://creativecommons.org/licenses/by/4.0/, as recorded in the arXiv licence metadata for this exact version and shown on the abstract page https://arxiv.org/abs/2604.08840v1. Full licence notice: \texttt{licenses/arxiv-2604.08840-CC-BY-4.0.txt}.\par\smallskip

\noindent\textit{Editorial scope.} Counted unit: the article's theorem theorem:equilibria (source0.tex lines 420--431). The article's complete original proof of it is delivered with it (source0.tex lines 432--450, one \texttt{proof} environment of 19 lines). No mathematics is written, repaired or re-derived: the statement and the proof are the article's own lines, in the article's own notation, and no retained line is substituted or re-flowed.  Retained uncounted as the source-local context the proof needs: lines 239--242; lines 325--327; lines 352--354.  Nothing was omitted from any retained span, so the package carries no omission record.  No retained line contains a citation, so the wrapper carries no bibliography and no bibliography entry.  No retained line contains a figure environment, an image-inclusion command or drawing code, so no image is delivered and the package declares no asset dependency.  Editorial formatting only, in addition: an AMS \texttt{amsart} preamble that restates the article's own declarations and macros used by the retained lines, namely the environments \texttt{assumption}, \texttt{proposition}, \texttt{theorem} on separate counters, together with the standard packages those lines need.  The retained environments print in this document as Assumption 1, Proposition 1, Theorem 1 in order of appearance, whereas the article prints the same items with the numbering of its own full document.  The complete native source of the article as distributed in the arXiv e-print for this exact version is bundled unchanged as \texttt{sources/198176/source0.tex}.
\begin{align}\label{eq:discriminant}
    \delta_i(z_i,\vect{z}_{-i}) &=  \alpha_i\left( \dfrac{r}{n}-1 \right) \\ \notag
    &+ \frac{\beta_i\lambda_i}{\beta_i+\lambda_i}\bigg(\gamma_iu_i+(1-\gamma_i)\sum_{j\in\mathcal{V}}w_{ij}y_j-\frac{1}{2}\bigg).
\end{align} 

\begin{assumption}\label{a:0}
    All players have zero prejudice, i.e. $\gamma_i = 0, \forall i \in \mathcal{V}$, and $\alpha_i, \beta_i, \lambda_i \in (0,1), \forall i \in \mathcal{V}$.
\end{assumption}

\begin{proposition}\label{prop:consensus}
    Consider the coevolutionary model under Assumption~\ref{a:0}, and let $\vect{z}^*=(x^*,y^*)$ be an equilibrium of the model. There holds $\vect x^* = \vect 0 \Rightarrow \vect y^* = \vect 0$ and $\vect x^* = \vect 1 \Rightarrow \vect y^* = \vect 1$. That is, an equilibrium $\vect z^*$ with an action consensus must in fact also have an opinion consensus.  
\end{proposition}

\begin{theorem}\label{theorem:equilibria}
    Under Assumption~\ref{a:0}, the all-defection consensus state $(\vect{0},\vect{0})$ is always an equilibrium of the coevolutionary model. Moreover, it is the unique equilibrium of the coevolutionary model if there holds
    \begin{align}\label{eq:all_D_unique}
        \frac{\beta_i\lambda_i}{\beta_i+\lambda_i} \leq 2\alpha_i\left(1-\frac{r}{n}\right),
    \end{align}
    for all $i\in\mathcal V$. If instead
       \begin{align}\label{eq:all_coop}
        \frac{\beta_i\lambda_i}{\beta_i+\lambda_i} > 2\alpha_i\left(1-\frac{r}{n}\right),
    \end{align}
    %\eqref{eq:all_D_unique} does not hold 
    for all $i\in\mathcal V$, then the all-cooperation consensus state $(\vect{1},\vect{1})$ is an equilibrium.
\end{theorem}

\begin{proof}
    To prove the first statement, we need to show that at an equilibrium $\vect z^* = (\vect x^*,\vect y^*)$, it is always true that i) $\vect x^* = \vect 0 \Rightarrow \vect y^* = \vect 0$, and ii) $\vect y^* = \vect 0 \Rightarrow \vect x^* = \vect 0$. The first part is immediate from Proposition~\ref{prop:consensus}. For the second part, observe that when $\vect y^* = \vect 0$, \eqref{eq:discriminant} evaluates to be
    \begin{equation}
        \delta_i(\vect y^*) = \alpha_i\bigg(\frac{r}{n}-1\bigg) - \frac{\beta_i\lambda_i}{2(\beta_i+\lambda_i)} \leq 0
    \end{equation}  
    because $r<n$ (recall $r$ is the public good multiplier). In other words, $x_i^* = 0$ for all $i$, hence $\vect z^* = (\vect 0,\vect 0)$ is always an equilibrium of the coevolutionary model.

    To prove the second statement, notice that if
    \begin{equation}\label{eq:01}
        \dfrac{\alpha_i(1-\alpha_i)\Big(1-\dfrac{r}{n}\Big)}{\beta_i\lambda_i}+\frac{1}{2} \geq1,
    \end{equation}
    then $\delta_i(\vect y^*) \leq 0$ holds independent of the particular value of $\vect y^*$, making $\vect z^*=(\vect 0,\vect 0)$ the unique equilibrium. Rearranging \eqref{eq:01} yields \eqref{eq:all_D_unique}. 

    To prove the all-cooperation consensus state is an equilibrium we must show that at an equilibrium $\vect z^* = (\vect x^*,\vect y^*)$, i) $\vect x^* = \vect 1 \Rightarrow \vect y^* = \vect 1$, and ii) $\vect y^* = \vect 1 \Rightarrow \vect x^* = \vect 1$. The first part is obtained from Proposition~\ref{prop:consensus}. For the second part, observe that under the assumption $\vect y^* = \vect 1$, \eqref{eq:discriminant} evaluates to be
    \begin{equation}\label{eq:02}
        \delta_i(\vect y^*) = \alpha_i\bigg(\frac{r}{n}-1\bigg) + \frac{\beta_i\lambda_i}{2(\beta_i+\lambda_i)}.
    \end{equation} 
    For a player to cooperate, $\delta_i(\vect y^*)>0$. Setting the condition in \eqref{eq:02} to be greater than zero and rearranging yields \eqref{eq:all_coop}. Therefore, if \eqref{eq:all_coop} holds for all $i \in \mathcal{V}$ then the all-cooperation consensus state is an equilibrium. %The condition in \eqref{eq:all_D_unique} shows the trade-off between self-consistency and opinion dynamics (left) and the public goods game (right).
\end{proof}

\end{document}
